\section{AC--MOT V1: Quality-Oriented Optimization}
An early heuristic prototype was used to identify the controllable detector
variables and the detector--tracker causal interface. We define V1 as the
first formal optimization-derived AC--MOT configuration evaluated under the
frozen experimental protocol. The prototype remains fully documented in
Online Resource~1, but it is not treated as the primary AC--MOT generation.

\subsection{Optimization design}
V1 fixed pretrained YOLOv8n and a tuned ByteTrack profile and searched only on
the VisDrone2019-MOT validation split. The test split was not accessed during
search. At analysis step $j$, the controller formed five inexpensive cues:
crowd count, the fraction of prior boxes below $32\times32$ pixels, edge
density, low brightness, and low Laplacian variance. Each cue was calibrated
empirically from fixed-detector validation outputs or validation image
statistics, and the SCI was their weighted sum,
\begin{equation}
\begin{aligned}
S_j&=\sum_{k\in\{c,t,e,n,b\}}w_k x_{k,j},\\
w_k&\geq0,\quad\sum_k w_k=1.
\end{aligned}
\end{equation}
The seven-entry moving average was evaluated every ten frames. Confidence and
NMS were linearly interpolated between the selected SCI=0 and SCI=1 endpoints,
while two sorted SCI thresholds selected one of the three resolution levels.

The canonical Optuna/TPE study jointly searched five normalized cue weights,
the confidence endpoints, the NMS endpoints, and the two SCI thresholds
\cite{akiba2019optuna}. The resolution levels $512,928,960$, supported
confidence values $\{0.25,0.30,0.35,0.40,0.45\}$, supported NMS values from
$0.30$ through $0.70$, and the temporal settings $W=7$ and stride $10$ were
provided by preceding validation screens or the frozen temporal design; they
were not falsely presented as variables of the final joint search.

The study recorded the multi-objective directions maximize MOTA, minimize IDS,
and maximize FPS. The final validation-only feasibility rule required the FPS
gate and $\mathrm{IDS}\leq271$ (the OLD--A3 validation reference), after which
the highest MOTA was selected, with lower IDS, HOTA, IDF1, and FPS as tie
breaks. Trial24 was frozen from this procedure. Its exact controller values
were
\begin{equation*}
\begin{split}
&w_c=0.1294928,\quad w_t=0.2217477,\\
&w_e=0.4337134,\quad w_n=0.0535577,\\
&w_b=0.1614885,\\
&(c_{easy},c_{hard})=(0.30,0.40),\\
&(n_{easy},n_{hard})=(0.35,0.35),\\
&(\tau_{mid},\tau_{high})=(0.1353494,0.2872868).
\end{split}
\end{equation*}

\subsection{Historical held-out profile}
On the historical held-out comparison, the baseline obtained MOTA/HOTA/IDF1
$19.729/28.430/32.724$, IDS $1235$, and 36.528 FPS. V1 obtained
$26.948/33.835/41.546$, IDS $1184$, and 38.985 FPS. The archived paired
bootstrap supports V1--baseline gains of $+7.2195$ MOTA, $+5.4051$ HOTA, and
$+8.8220$ IDF1, with 95\% intervals $[5.4362,9.2934]$,
$[4.1273,6.9022]$, and $[6.9005,11.0537]$, respectively. The IDS reduction
of 51 has interval $[-114,219]$ and is therefore not a statistically
established identity improvement. V1 is a quality-oriented profile: it was
not selected as the lowest-IDS profile and it is not claimed to be a universal
operating point.
